\documentclass[11pt]{article}
\usepackage[margin=1in]{geometry}
\usepackage[T1]{fontenc}
\usepackage[utf8]{inputenc}
\usepackage{amsmath,amssymb,amsthm}
\usepackage{enumitem}

\title{ICPC World Finals 2021\\C. Fair Division}
\author{}
\date{}

\begin{document}
\maketitle

\section*{Problem Summary}

There are $n$ pirates in a circle. Every time a pirate receives some loot, they keep a fraction
$f \in (0,1)$ and pass the remaining fraction $1-f$ to the next pirate. This continues forever.

We must find a rational fraction $f = \frac{p}{q}$ such that the \emph{total} amount eventually assigned to
each pirate is an integer. Among all valid fractions, we want the smallest denominator $q$, and among those
the smallest numerator $p$.

\section*{From Infinite Process to Integer Geometric Progression}

Let
\[
    r = 1 - f.
\]
If the total loot is $m$, then pirate $i$ receives
\[
    m \cdot f \cdot r^i
\]
in the first round, then the same amount multiplied by $r^n$ in the next round, and so on. Therefore the
total share of pirate $i$ is
\[
    x_i = m \cdot \frac{f r^i}{1-r^n}.
\]

So the pirates' final shares form a geometric progression:
\[
    x_{i+1} = r x_i.
\]
The problem is equivalent to asking for \emph{integer} values
\[
    x_0, x_1, \dots, x_{n-1}
\]
with common ratio $r \in (0,1)$ and total sum $m$.

\section*{Lowest-Term Form}

Write
\[
    r = \frac{a}{q}
\]
in lowest terms, where $1 \le a < q$ and $\gcd(a,q)=1$. Then every share has the form
\[
    x_i = y \cdot q^{n-1-i} a^i
\]
for some positive integer $y$. Indeed, the denominators force the first term to contain a factor
$q^{n-1}$, and after that all terms are integral automatically.

Summing the progression gives
\[
    m
    = y \sum_{i=0}^{n-1} q^{n-1-i} a^i
    = y \cdot D(q,a),
\]
where
\[
    D(q,a) = q^{n-1} + q^{n-2}a + \cdots + a^{n-1}.
\]

Therefore a fraction is valid exactly when
\[
    D(q,a) \mid m.
\]

Since
\[
    f = 1-r = \frac{q-a}{q},
\]
the output numerator is
\[
    p = q-a.
\]
Because $\gcd(a,q)=1$, we also have $\gcd(p,q)=1$, so this denominator is already minimal for that choice.

\section*{Search Strategy}

We search denominators $q=2,3,4,\dots$ in increasing order. For a fixed $q$, we try numerators
$p=1,2,\dots,q-1$ in increasing order, which is exactly the required tie-break order.

For one candidate:
\begin{itemize}[leftmargin=*]
    \item the remainder ratio is $a = q-p$;
    \item the candidate works iff $\gcd(p,q)=1$ and $D(q,a)$ divides $m$.
\end{itemize}

The smallest possible value of $D(q,a)$ for this denominator happens when $a=1$, namely
\[
    1 + q + q^2 + \cdots + q^{n-1}.
\]
Once even this minimum exceeds $m$, no larger denominator can work and the search stops.

\section*{Useful Bound}

The absolute smallest divisor over \emph{all} valid fractions occurs at $q=2$ and $a=1$, which gives
\[
    2^n - 1.
\]
If $n \ge 60$, then
\[
    2^n - 1 > 10^{18} \ge m,
\]
so no solution exists at all. This removes the large-$n$ cases immediately.

\section*{Algorithm}

\begin{enumerate}[leftmargin=*]
    \item If $n \ge 60$, output \texttt{impossible}.
    \item For $q = 2,3,\dots$:
    \begin{itemize}[leftmargin=*]
        \item compute $D(q,1)$; if it already exceeds $m$, stop and output \texttt{impossible};
        \item for each $p=1,2,\dots,q-1$:
        \begin{itemize}[leftmargin=*]
            \item skip it if $\gcd(p,q) \ne 1$;
            \item let $a=q-p$;
            \item compute $D(q,a)$ with overflow capping at $m+1$;
            \item if $D(q,a)$ divides $m$, output $p$ and $q$.
        \end{itemize}
    \end{itemize}
\end{enumerate}

To evaluate
\[
    D(q,a) = q^{n-1} + q^{n-2}a + \cdots + a^{n-1},
\]
we use the recurrence
\[
    S_1 = 1, \qquad S_{k+1} = qS_k + a^k.
\]
This avoids separate exponentiation.

\section*{Correctness Proof}

We prove that the algorithm returns the required fraction whenever one exists.

\paragraph{Lemma 1.}
If the final shares are integers, then they form an integer geometric progression with ratio $r=1-f$.

\paragraph{Proof.}
Pirate $i$ receives a fraction $f$ of the amount arriving at their position in every round. Between pirate
$i$ and pirate $i+1$, the remaining loot is multiplied by $r=1-f$. Therefore every contribution to pirate
$i+1$ is exactly $r$ times the corresponding contribution to pirate $i$, and summing over all rounds gives
$x_{i+1}=r x_i$. \qed

\paragraph{Lemma 2.}
Let $r=\frac{a}{q}$ be in lowest terms. There exists a valid division if and only if
$D(q,a)$ divides $m$.

\paragraph{Proof.}
Suppose a valid division exists. By Lemma 1, the shares are an integer geometric progression with ratio
$\frac{a}{q}$. Because $\gcd(a,q)=1$, integrality forces the first term to be a multiple of $q^{n-1}$, so
the progression has the form
\[
    y q^{n-1},\; y q^{n-2}a,\; \dots,\; y a^{n-1}
\]
for some positive integer $y$. Summing yields $m = y D(q,a)$, so $D(q,a)\mid m$.

Conversely, if $D(q,a)\mid m$, let $y = \frac{m}{D(q,a)}$ and define
\[
    x_i = y q^{n-1-i} a^i.
\]
These are integers, their ratio is $\frac{a}{q}$, and their sum is exactly $m$. Taking
$f = 1-\frac{a}{q} = \frac{q-a}{q}$ produces exactly those total shares. \qed

\paragraph{Lemma 3.}
For a fixed denominator $q$, checking $p=1,2,\dots,q-1$ in increasing order finds the correct tie-break
winner for that denominator.

\paragraph{Proof.}
The output fraction is $f=\frac{p}{q}$, and the statement asks for the smallest numerator among equal
denominators. The loop checks these numerators in exactly that order and stops at the first valid one. \qed

\paragraph{Lemma 4.}
If $D(q,1) > m$, then no denominator larger than $q$ can be valid.

\paragraph{Proof.}
For fixed $q$, the smallest possible remainder numerator is $a=1$, which minimizes $D(q,a)$. Also
$D(q,1)$ increases with $q$. Therefore if the minimum value for this denominator already exceeds $m$, then
every candidate with denominator at least $q$ has divisor larger than $m$ and cannot divide $m$. \qed

\paragraph{Theorem.}
The algorithm outputs \texttt{impossible} exactly when no suitable fraction exists; otherwise it outputs the
valid fraction with smallest denominator, and among those, smallest numerator.

\paragraph{Proof.}
By Lemma 2, a candidate fraction is valid exactly when the algorithm's divisibility test succeeds. By
Lemma 4, once the outer loop stops, no unseen denominator can work. Since denominators are tested in
increasing order and, for each denominator, numerators are tested in increasing order, the first valid
fraction found is exactly the one required by the problem. If no candidate is found before the stopping
condition, then no valid fraction exists. \qed

\section*{Complexity Analysis}

If $n \ge 60$, the answer is immediate.

Otherwise, every candidate evaluation takes $O(n)$ time, and the outer search stops quickly because
$D(q,1)$ grows as a degree-$(n-1)$ polynomial in $q$. The memory usage is $O(1)$.

\section*{Implementation Notes}

\begin{itemize}[leftmargin=*]
    \item The recurrence for $D(q,a)$ is evaluated with \texttt{\_\_int128} and capped at $m+1$ to avoid
    overflow.
    \item Because $\gcd(q-a,q)=\gcd(a,q)$, checking $\gcd(p,q)=1$ is enough.
\end{itemize}

\end{document}
